\section{Related Work}
\label{sec:related_work}

围绕乳腺组织病理图像（breast histopathology image）分析的相关工作，可以归纳为三条彼此关联的研究线索：数字病理（digital pathology）从人工设计特征（hand-crafted features）转向深度学习（deep learning）的总体演进，IDC 图像块级（patch-level）分类中端到端学习（end-to-end learning）与迁移学习（transfer learning）的应用，以及医学人工智能（medical AI）中围绕评估协议、数据泄漏（data leakage）与可复现性（reproducibility）的持续讨论。本文的工作并不试图在网络结构上继续堆叠复杂模块，而是立足于这三条线索的交汇处：以端到端深度学习基线为对象，在严格协议与统一形式化约束下开展可复验的实证研究。

\subsection{Deep Learning in Digital Pathology}

早期组织病理图像分析多依赖颜色统计、纹理描述子、形态学特征及其与浅层分类器的组合。这类方法具有一定可解释性，但对预处理质量、特征设计经验和场景先验高度敏感，难以稳健应对染色变异（stain variation）、组织异质性（tissue heterogeneity）和跨中心分布偏移（domain shift）\cite{Komura2018,Litjens2017}。随着卷积神经网络（convolutional neural network, CNN）在视觉任务中的成功，数字病理逐步转向由模型直接从像素中学习层级化表征（hierarchical representations）的范式\cite{Litjens2017,vanDerLaak2021}。

这一转向的意义并不只是“分类器更强”，而在于表示学习（representation learning）与最终诊断目标之间建立了更直接的耦合。对于病理图像而言，判别线索常常横跨细胞级别的核异型性（nuclear atypia）、腺体结构与局部组织排列等多个尺度，端到端模型能够在统一损失下共同吸收这些统计模式，而不必将其拆解为若干彼此割裂的人工模块\cite{vanDerLaak2021,Huang2025,Rad2025}。因此，深度学习已成为计算病理学中的主流方法基础，而“如何在特定任务中建立可信基线”逐渐取代“是否使用深度学习”本身，成为更关键的问题。

\subsection{End-to-End Learning for Breast Histopathology and IDC Classification}

在乳腺组织病理分析中，深度学习已经被广泛用于良恶性分类（benign--malignant classification）、亚型识别、分级（grading）以及更高层级的切片预测任务\cite{Araujo2017,Hameed2020,Yan2020}。Ara{\'u}jo 等较早展示了 CNN 在乳腺癌组织学图像分类中的有效性，说明深层表示可以显著替代传统特征工程\cite{Araujo2017}；随后，预训练骨干网络、注意力机制和更复杂的特征聚合策略被不断引入乳腺病理分类，以增强对细粒度纹理与局部上下文的建模能力\cite{Yan2020,Zou2022}。

具体到浸润性导管癌（invasive ductal carcinoma, IDC）识别，相关研究大致分为两类：一类面向全视野切片（whole-slide image, WSI）或局部区域的检测与定位，典型如 Cruz-Roa 等将 CNN 用于乳腺癌切片中的 IDC 区域检测\cite{CruzRoa2014}；另一类则围绕公开 patch 数据集开展局部二分类研究，即直接判断图像块是否包含 IDC。后者因算力开销较低、数据组织清晰而成为最常见的基准设定。已有研究表明，基于 ImageNet 预训练的 ResNet、DenseNet 等模型能够在这一任务上建立有竞争力的基线\cite{Celik2020,Abdolahi2020}。进一步地，Voon 等对多种 CNN 架构的比较显示，迁移学习后的端到端模型通常可以提供稳定性能，但网络更深或参数更多并不必然带来线性提升\cite{Voon2022}。这意味着，对于 IDC patch classification 而言，模型复杂度并不是唯一值得追逐的方向，实验边界与比较方式同样会决定结论的含义。

\subsection{Evaluation Protocols, Data Leakage, and Reproducibility}

与自然图像分类相比，医学影像数据更容易违反独立同分布（independent and identically distributed, i.i.d.）假设。一个患者往往对应多个切片、多个区域或大量高度相关的 patch；若划分时不尊重患者、切片或体数据（volume）的层级结构，训练集与验证集就可能共享强相关样本，从而使模型利用不应出现在泛化评估中的捷径信号（shortcut signals）。Yagis 等在脑 MRI 任务中展示了受试者级相关性被破坏时性能估计会发生显著偏移\cite{Yagis2021}；Tampu 等进一步证明，不恰当的切分策略足以显著膨胀公开 OCT 数据集上的测试性能\cite{Tampu2022}。虽然这些结果来自不同模态，但其方法学含义对病理 patch 数据同样成立。

对数字病理而言，这一问题往往更加隐蔽，因为 patch 数据通常源于对少量 WSI 的密集裁剪，相邻样本之间天然共享组织纹理、染色风格与扫描噪声。若直接采用 patch-wise 随机划分，模型就可能在训练阶段接触与验证样本高度相似的患者内模式，从而夸大其对未见患者的泛化能力。近年来的综述也不断强调，外部泛化不足、验证设计薄弱和结果难以复现，仍是数字病理走向可信临床应用的核心障碍\cite{vanDerLaak2021,McGenity2024,Rad2025}。

与协议问题相伴的，是可复现性对实验结论的持续影响。对于医学深度学习研究，最终性能不仅受网络结构影响，也受配置管理、随机种子、数据采样、损失加权、阈值选择、日志记录与 checkpoint 保存等工程因素影响。若这些变量缺乏系统控制，即使使用同一公开数据集，不同研究之间的结果也很难公平比较\cite{McGenity2024,Huang2025}。因此，在这一研究脉络下，patient-wise 划分、配置驱动实验和训练/评估解耦不应被视为外围实现细节，而应被视为方法有效性的组成条件。

\subsection{Positioning of This Work}

基于上述文献脉络，本文的定位可以概括为以下三个相互支撑的方面：
\begin{enumerate}
    \item 在实证层面，本文以可切换的预训练 ResNet 主干网络（backbone）和轻量分类头为统一起点，构建模块化、配置驱动且可复现的实验框架，并在统一条件下组织 backbone 基线比较与切分协议比较；
    \item 在协议层面，本文将 patient-wise split、类别不平衡控制、weighted sampler 与可追踪的数据准备流水线纳入同一数据流治理视角，以降低患者内相关性、采样偏置与训练分布失衡对结论的干扰；
    \item 在理论层面，本文进一步将 patient-grouped distribution、代价敏感训练（cost-sensitive training）与阈值化决策（thresholded decision rule）置于统一形式化表述之下，为经验结果提供可解释的统计学习参照。
\end{enumerate}

因此，本文与已有工作的区别并不主要体现在额外设计一个更复杂的分类头，而体现在：我们试图把 IDC patch classification 从一个容易被实验捷径误导的局部高分问题，重新组织为一个由实证基线、协议边界与形式分析共同支撑的研究对象。这样的定位既回应了端到端深度学习基线研究的需求，也回应了医学 AI 对严格评估与可复现性的要求，并为后续的方法比较、消融分析与系统扩展提供更可信的起点。
